% !TeX root = ../papers/nc-p3.tex
\section{吸引ファイバーの全周期と重み付き座標での判別}
\label{sec:torus-fiber-periods}
\begin{theorem}\label{thm:torus-fiber-periods}
定理\ref{thm:finite-torus-slices}および定理\ref{thm:torus-component-crossing}の
非閉軌道の全てについて、最小周期は次の通りである。
\begin{enumerate}
\item 単重み曲面の一般点 $A$ に落ちるものは全て4、単重み曲線の一般点 $B$ に落ちるものは全て2である。
\item 共通点 $w_*$ の重み付き射影空間上で、重み3と4の各二座標のうち一つだけを符号反転し、
残りを固定する対合を $F$ とする。$F$ の不動点がちょうど周期1、それ以外は2である。
\item 中央二重型の一般点 $C$ に落ちる空間の重み $1,3$ の各二座標を、符号対合 $D$ の
正負の固有順に取る。$F$ はこの二座標を交換する。重み2の三座標では
$D=\operatorname{diag}(1,1,-1)$, $F=\operatorname{diag}(1,-1,1)$ とする。
$F$ が $D$ 商上で固定する点がちょうど周期1、それ以外は2である。
\end{enumerate}
具体的に重み $j$ のベクトルを $v_j$ とすると、最後の場合はある $t\in k^\times$ と
$e\in\{0,1\}$ が全 $j$ で $Fv_j=t^j D^e v_j$ を満たすこと、二番目なら同じ式で $e=0$ とすることが必要十分である。
従って座標が消える部分や二つの符号成分が同時に現れる場合も含む。
\end{theorem}
\begin{proof}
位置巡回を基底群に加えた有限拡大は還元的である。閉軌道を保つ巡回と基底変更の組にも同変に
前二節のスライスを取れる。特異な共通点でも滑らかな周囲のスライスを同変に取るので、
零ファイバー上は実際に正負の線形重み空間の作用になる。群の有限拡大によって
元の同型関係を変えることはせず、巡回後に元の基底群軌道へ戻るかを計算する。
$A$ は既に閉極限の周期で決まる。

$B$ の一移動は重みを反転する同型でのみ閉点に戻せ、正負の二コピーを交換する。
二移動は正負を保ち、位置0に $\operatorname{diag}(b,1,1,1)$ を作用させた表示では
単に二つの反対位置を同時交換する位数2の作用である。
$b=2$ でその正側の軌道法表現への作用を計算すると全て恒等となる。負側も一移動による共役で同じである。

共通点では向きを保つ一移動がある。四テンソルの座標を一つ巡回し、新位置2の数字0を $1/2$ 倍、
新位置3の数字0を2倍すれば固定される（他は恒等）。
正側での各重みの作用は全て二乗が恒等で、負固有空間の次元が重み順に $0,0,1,1,0$ である。
従って二移動は全点で恒等、一移動の不動条件は記載したものとなる。

$C$ では巡回と二中間文字の交換を組み合わせれば向きを保つ。
$\alpha=2$ の基底で、この単項式変換に対し四位置に必要な対角倍率を作用させた写像を $L$ とすると、
重み $j$ 上で $L^2=(\sqrt{2}/4)^j I$, $LD=(-1)^j DL$ となる。
トーラスの元を掛けて全重みで $F^2=1$ にできる。奇数重みの二座標は
$D$ の固有空間が一つずつで互いに交換されるので、基底を選んで表示通りとなる。
偶数重みでは $D$ の正側で $F$ は正負一つずつ、負側で正一つとなる。
持上げに使ったトーラス倍率と $D$ は既に同型関係に含まれる。

\path{computations/verify_torus_fiber_periods.py} は元の窓上で各変換を検算し、
二切断の微分の核から軌道方向を除いた各空間でこの作用を計算する。
各族の稠密開では既に重み別法束が局所自由であり、シフトの正規化に必要な根をエタールに添加すれば
有限位数作用の固有束の階数は一定である。実際、中央二重型で巡回の四乗は向きと中間二線を保ち、
二移動自体は両対称表示の二交換による位数2の持上げを持つ。
より直接には一移動の四乗はトーラスと $D$ の積であり、トーラスの四乗根で除いた後の位数は高々8である。
これで一般点にも標本と同じ零ファイバー判定が移る。
\end{proof}
